\begin{frame}{Operadores relacionais: comparar valores}
\begin{center}
\begin{tabular}{lll}
\toprule
Operador & Significado & Exemplo verdadeiro \\
\midrule
\texttt{<} & Menor que & \texttt{20 < 30} \\
\texttt{<=} & Menor ou igual & \texttt{30 <= 30} \\
\texttt{>} & Maior que & \texttt{40 > 30} \\
\texttt{>=} & Maior ou igual & \texttt{30 >= 30} \\
\texttt{==} & Igual a & \texttt{5 == 5} \\
\texttt{!=} & Diferente de & \texttt{5 != 4} \\
\bottomrule
\end{tabular}
\end{center}
Em C, comparações produzem \texttt{int}: 1 para verdadeiro e 0 para falso. Podemos armazenar o resultado em \texttt{bool}.
\end{frame}
